%% kapa: exam=1 round=34 year=2023 period=1 form=A subject=civil_law subject_ko=민법 seq=23 number=23
%% source: Q-Net 감정평가사 기출문제 (cst00309) artlSeq=5211789 file=raw/1008/5211789/2023년도 제34회 감정평가사 1차시험 1교시.pdf
%% conversion: kordoc 4.24.51 -> tools/gen1.py (auto); review=proofread-sonnet; flags=none
\begin{question}{23}
등기의 유효요건에 관한 설명으로 옳지 않은 것은? (다툼이 있으면 판례에 따름)
\begin{choices}
  \item 물권에 관한 등기가 원인 없이 말소되더라도 특별한 사정이 없는 한 그 물권의 효력에는 영향을 미치지 않는다.
  \item 미등기건물의 승계취득자가 원시취득자와의 합의에 따라 직접 소유권보존등기를 마친 경우, 그 등기는 실체관계에 부합하는 등기로서 유효하다.
  \item 멸실된 건물의 보존등기를 멸실 후에 신축된 건물의 보존등기로 유용할 수 없다.
  \item 중복된 소유권보존등기의 등기명의인이 동일인이 아닌 경우, 선등기가 원인무효가 아닌 한 후등기는 무효이다.
  \item 토지거래허가구역 내의 토지에 대한 최초매도인과 최후매수인 사이의 중간생략등기에 관한 합의만 있더라도, 그에 따라 이루어진 중간생략등기는 실체관계에 부합하는 등기로서 유효하다.
\end{choices}
\end{question}
\answer{5}
